\section{Conclusion and Future Directions}
\label{sec:conclusion}

In this work, we challenged the traditional coordinate-dependent spatial paradigm of model merging and introduced \textbf{SpectralMerge}, a novel framework that reformulates parameter consolidation in the frequency domain. By applying the Discrete Cosine Transform (DCT-II) to layer-wise merging coefficients, we demonstrated that the underlying parameter sensitivity can be elegantly regularized and analyzed in spectral coordinates. We proposed two frequency-domain formulations: \textbf{SpectralMerge-LP} (Low-Pass hard filtering) and \textbf{SpectralMerge-Reg} (quadratic soft spectral decay). 

Our extensive empirical evaluations across 30 random seeds on Vision Transformer (ViT-B/32) multi-task landscapes have yielded several major findings. First, SpectralMerge outpeforms all competitive baselines on standard clean streams under test-time adaptation and offline few-shot validation tuning. Second, we systematically refuted the Overfitting-Optimizer Paradox under extreme sample complexity ($M=5$), showing that our spectral constraints act as robust analytical low-pass filters that completely reject high-frequency transductive noise. Third, we demonstrated that SpectralMerge exhibits exceptional resilience and graceful degradation under severe validation target selection bias and adversarial stream noise.

From a broader perspective, we believe this work merely scratches the surface of spectral representations in deep learning consolidation. By showing that orthogonal frequency bases offer superior numerical conditioning and structural regularization compared to spatial or polynomial coordinates, we hope to inspire a new paradigm of frequency-domain parameter management. 

Exciting future directions include:
\begin{itemize}
    \item \textbf{Scaling to Pre-trained Foundation Models:} While we demonstrated SpectralMerge's effectiveness on continuous simulation and a heterogeneous physical multi-layer perceptron, applying it to massive pre-trained foundation models (such as merging fine-tuned \texttt{RoBERTa-base}, \texttt{ViT-B/16}, or Large Language Models like \texttt{Llama-3}) represents a highly promising avenue. In practice, this would involve extracting layer-wise task vectors for deep networks, applying the 1D DCT-II to partition homogeneous attention/feedforward projections, and using standard PEFT libraries or zero-shot datasets to optimize the low-frequency spectral coefficients, serving as a highly scalable alternative for massive-scale pre-trained model consolidation.
    \item \textbf{PEFT-Induced Discontinuities and Hybrid Spectral Designs:} Merging models trained via Parameter-Efficient Fine-Tuning (such as LoRA or prefix-tuning) or localized adaptation (where only specific blocks or deep layers are fine-tuned, as in our ResNet-18 evaluations) introduces a sharp step-function discontinuity in task-vector magnitudes across depth. Because step functions exhibit infinite frequency support, rigid hard cutoffs (such as SpectralMerge-LP and LP-Adaptive) catastrophically underfit by filtering out the high-frequency components required to represent these localized active layers. While soft regularization (SpectralMerge-Reg) successfully overcomes this by allowing gradients to activate localized high-frequency coordinates, future research should develop hybrid spectral designs. These could include piece-wise localized DCT transforms or adaptive, layer-masked frequency windows to explicitly capture PEFT-induced boundary transitions.
    \item \textbf{Synergy with Sign-and-Magnitude Filtering:} Integrating SpectralMerge with orthogonal parameter-space filtering techniques (such as TIES-Merging's sign consensus and magnitude pruning or DARE's delta-weight dropping). In this hybrid paradigm, sign conflict resolution and sparsification are first executed to isolate core task updates, and SpectralMerge is subsequently applied to optimize and regularize the layer-wise scaling of the filtered parameters in the frequency domain. This co-design holds immense potential for non-interfering parameter consolidation.
    \item \textbf{Adaptive and Dynamic Spectral Bandwidth:} In this work, we took a key step toward this direction by implementing and validating an \textbf{Adaptive Bandwidth SpectralMerge (LP-Adaptive)} formulation on physical PyTorch networks (Section~\ref{sec:experiments}). We showed that starting optimization with highly regularized low-frequency components ($F=1$) and dynamically expanding to higher-frequency capacities ($F=5$) improves performance over fixed cutoffs. A natural extension is to formulate a fully continuous, data-driven bandwidth mechanism where the cutoff is governed dynamically by an adaptive threshold on the running validation loss or energy spectral density, allowing the model to balance regularization and representation capacity completely on the fly.
    \item \textbf{Multidimensional and Joint Spectral Merging in Large-Scale Pools ($K \ge 8$):} Extending the 1D DCT-II representation across network depth to 2D or 3D spectral transforms is especially crucial as we scale to larger multi-task configurations (e.g., $K \ge 8$ or $K \ge 12$ task experts, typical in modern LLM or visual adapter merging). As the task pool size $K$ grows, task interference and weight representation clashes escalate exponentially. Optimizing separate 1D depth trajectories for each task independently scales the parameter space to $K \times L$, which heightens the risk of overfitting under extreme data scarcity. 
    
    To overcome this, we propose treating the combining coefficient matrix $\alpha \in \mathbb{R}^{K \times L}$ as a unified 2D discrete parameter signal and applying a 2D DCT-II transform. This maps the spatial depth-and-task coordinates into joint 2D spectral coordinates. Because task experts fine-tuned from a shared initialization exhibit highly correlated scaling behaviors, the 2D DCT-II can capture these high-level inter-task orthogonal correlations alongside depth-wise continuity, packing over $95\%$ of the joint energy into a tiny corner of low-frequency 2D coordinates. This enables the model to compress the joint optimization search space drastically (e.g., to only $4 \times 3$ trainable parameters for $K=12$ tasks) while maintaining rich representation capacity, mitigating task interference, and guaranteeing robust multi-task generalization.
    \item \textbf{Computational Scaling to Fine-Grained (Parameter-wise) Merging:} In this work, we focus on layer-wise or block-wise merging, where the number of parameters $L \le 96$ is small, and the 1D DCT-II computation is virtually instantaneous. However, if one wishes to scale SpectralMerge to fine-grained, parameter-wise merging coefficients (where each individual weight in a network has its own trainable combining coefficient), performing the forward and inverse transforms on millions of independent parameters would become a significant computational bottleneck. To overcome this limitation, future work could design batch-parallelized, GPU-efficient 1D/2D DCT and IDCT algorithms implemented via custom CUDA kernels or leveraging PyTorch's native Fast Fourier Transforms (FFT) with discrete cosine symmetry extensions. This would enable backpropagating through high-dimensional spectral transforms in parallel, allowing fine-grained parameter-wise merging to scale efficiently to massive architectures. 

    Furthermore, we identify a fascinating open question regarding parameter-wise sensitivity profiles: while global layer-wise scaling trajectories are naturally smooth, individual parameter-wise sensitivities are likely to exhibit highly localized, high-frequency components across layers. Consequently, a global cosine transform (DCT) might suffer from spectral leakage when representing these sharp, localized variations. To resolve this, future research should explore localized spectral bases, such as multi-resolution \textbf{Discrete Wavelet Transforms (DWT)}, to capture localized parameter-wise high-frequency dynamics while maintaining frequency-domain regularization and compression. 

    Specifically, let $h[n]$ and $g[n]$ represent the low-pass (scaling) and high-pass (wavelet) analysis filters of a standard orthogonal wavelet family (such as Haar or Daubechies wavelets). At each decomposition level $j \in \{1, \dots, J\}$, the spatial coefficient signal $\vec{\alpha}_k$ is recursively decomposed into scaling (approximation) coefficients $a_{k, j}[n]$ and detail (wavelet) coefficients $d_{k, j}[n]$ via convolution and downsampling:
    \begin{align}
        a_{k, j}[n] &= \sum_{m} h[m - 2n] a_{k, j-1}[m], \\
        d_{k, j}[n] &= \sum_{m} g[m - 2n] a_{k, j-1}[m],
    \end{align}
    where $a_{k,0}[m] = \alpha_k(m)$ denotes the original spatial coefficient at layer $m$. For the Haar wavelet family, these filters are defined as $h = [1/\sqrt{2}, 1/\sqrt{2}]$ and $g = [1/\sqrt{2}, -1/\sqrt{2}]$. For Daubechies wavelets, the filter support is wider (e.g., Db4 has 4 filter coefficients) to guarantee higher-order vanishing moments, enabling the representation of more complex local polynomial trends. 

    The original spatial coefficient signal can be reconstructed perfectly via the Inverse Discrete Wavelet Transform (IDWT):
    \begin{align}
        a_{k, j-1}[m] &= \sum_{n} \tilde{h}[m - 2n] a_{k,j}[n] \nonumber \\
        &\quad + \sum_{n} \tilde{g}[m - 2n] d_{k,j}[n],
    \end{align}
    where $\tilde{h}$ and $\tilde{g}$ are the corresponding synthesis low-pass and high-pass filters. 

    In this multi-resolution framework, a wavelet-based regularization penalty can be formulated by penalizing the detail coefficients progressively across scales:
    \begin{equation}
        \mathcal{R}_{wavelet}(\Theta_{wavelet}) = \sum_{k=1}^K \sum_{j=1}^J 2^{2j} \|d_{k,j}\|_2^2
    \end{equation}
    where the penalty scale factor $2^{2j}$ scales exponentially with the decomposition level $j$ (reflecting increasingly finer spatial resolution scales). This wavelet decay penalty behaves as a highly localized structural regularizer: it allows global, coarse-scale trajectories (represented by the unpenalized approximation coefficients $a_{k, J}$) to adapt freely, while progressively suppressing high-frequency noise at finer resolution scales. Because wavelets possess compact spatial support, this formulation enables the model to capture sharp localized transitions (e.g., PEFT-induced step discontinuities) without suffering from the global spectral leakage (such as Gibbs-like oscillations) that would occur under a global cosine or Fourier basis, ensuring superior localized parameter consolidation.
    \item \textbf{Graph Spectral Fusion:} Applying Graph Fourier Transforms (GFT) to model the network connectome directly, enabling model merging on irregular non-Euclidean parameter topological backbones.
    \item \textbf{Spectral Denoising Hypernetworks:} Training generative weight-space diffusion models directly in the spectral domain to generate optimal multi-task parameters, bypassing spatial non-convex energy barriers.
\end{itemize}
We hope that SpectralMerge serves as a robust and mathematically elegant foundation for these future, paradigm-shifting research directions.
